% Source: https://artofproblemsolving.com/community/c3h2227851_perfectabundantdeficient_numbers
% Stable identity: aops-topic-2227851-q01
\textbf{(a)} Let $q=2^p-1$ be prime and $n=2^{p-1}q$. Since $q$ is odd, the divisors of $n$ are $2^i$ and $2^iq$ for $0\le i\le p-1$, so
\[\sigma(n)=\left(1+2+\cdots+2^{p-1}\right)(1+q)=\left(2^p-1\right)2^p=2\cdot2^{p-1}q=2n .\]
So $n$ is perfect (for example $6=2\cdot3$, $28=4\cdot7$, $496=16\cdot31$). $\blacksquare$

\textbf{(b)} The divisors of $q^k$ are $1,q,\ldots,q^k$. Since $q\ge2$,
\[1+q+\cdots+q^{k-1}=\dfrac{q^k-1}{q-1}\le q^k-1<q^k,\]
so $\sigma\left(q^k\right)=\left(1+q+\cdots+q^{k-1}\right)+q^k<2q^k$, and $q^k$ is deficient. $\blacksquare$

\textbf{(c)} Let $m=kn$ with $k\ge2$. For every divisor $d$ of $n$, the number $kd$ divides $m$, and these divisors add up to $k\sigma(n)=2kn=2m$. They are all at least $k\ge2$, so the divisor $1$ of $m$ is not among them. Hence
\[\sigma(m)\ge2m+1>2m,\]
and $m$ is abundant. Taking the perfect number $n=6$, every multiple $6k$ with $k\ge2$, that is $\boxed{12,18,24,30,\ldots}$, is abundant (for example $\sigma(12)=28>24$). $\blacksquare$
